\documentclass[11pt]{article}
\usepackage[margin=1in]{geometry}
\usepackage{amsmath,amssymb,amsthm,booktabs,longtable,hyperref,xurl}
\newcommand{\F}{\mathbb F}
\newcommand{\End}{\operatorname{End}}
\newcommand{\Hom}{\operatorname{Hom}}
\newcommand{\Aut}{\operatorname{Aut}}
\newcommand{\GL}{\operatorname{GL}}
\newcommand{\rank}{\operatorname{rank}}
\newcommand{\im}{\operatorname{im}}
\newcommand{\gb}[2]{{#1\brack #2}_q}
\title{Independent review of the acquired two-flag $E_6$ counting extension}
\author{Astra P3 continuation: independent foundations audit}
\date{28 October 2026}
\begin{document}\maketitle

\section*{Disposition and scope of this review}
\textbf{No required correction was found in the reviewed theorem or production
implementation.} The new assembly removes equality of the two nonzero target
carriers, while retaining the acquired two-factor regular matrix-dual and
zero-or-one joint-simple-copy target hypotheses. It admits every affine phase
under those hypotheses. Its subspace census is an explicit finite computation,
not an assumed orbit-counting oracle. The polynomial claim is uniform in field
degree and matrix sizes, with a large fixed exponent paid explicitly.

This review independently checks the argument, production implementation,
inherited phase orientation and native wrapper contract. It reads the saved
executed evidence and the audit implementation; it does not represent that
evidence as a fresh independent rerun. No additional execution was justified by
a remaining concrete risk, so no production or broad audit tests were rerun.
The earlier acquired universal-product fiber formula and raw frontend are
retained dependencies, not re-proved in full here. No recovered source or
production file was edited.

The principal reviewed files are
\path{new_work/phase_constructive/TWO_FACTOR_ARBITRARY_TARGET_PHASE_THEOREM.tex}
and \path{new_work/phase_constructive/e6_phase_counter.py}.
The exact SHA-256 identities reviewed were, respectively,
\begin{quote}\small\ttfamily
846b2bbf0f45e0cefb8e5b9897806470c43e478e00b8f701abd212ed9e26622a\\
de2ee0d65b8cf53c285e3760697357da6222b1723408ccf5b80f5f75c01021e1.
\end{quote}
The audit code and saved results had identities
\begin{quote}\small\ttfamily
b8560622dfb38b8e00b8645720d453631de97cd79c20b08b82d575817b23dea9\\
49a51560c720c137b87e65e4a97d5bcf37e0cb17062861bf2a5c8f8257d7d5a6.
\end{quote}

\section{Classification completeness and finite-field descent}
The quiver has vertices $(R,J,E_1,F_1,E_2,F_2)$ and arrows
$J\to R$, $E_i\to R$, $F_i\to E_i$. Its undirected graph has arms
of lengths $1,2,2$, hence is $E_6$. The production catalogue is the same
29-type catalogue described in the proof: eighteen ambient-one types, nine
ambient-two types, and two ambient-three types.

Each listed type is absolutely indecomposable, and this does not depend on
probabilistic field selection or on an experimental root calculation. In a
thin type, any endomorphism is fixed by its scalar ambient map. In a plane
type the two base arms force preservation of the lines $\langle e_1\rangle$
and $\langle e_2\rangle$, and $J$ forces preservation of
$\langle e_1+e_2\rangle$. A linear map preserving all three is scalar.
For each ambient-three type, the base flags preserve the three individual
lines $F_1=\langle e_1\rangle$, $E_1\cap E_2=\langle e_2\rangle$,
and $F_2=\langle e_3\rangle$. The remaining $J$ is either the diagonal
line or the plane with basis $e_1+e_2,e_2+e_3$. Its preservation forces
the three diagonal entries to coincide. These arguments hold in
characteristic two and after every scalar extension. Thus every listed
endomorphism ring is exactly the scalar field.

The 29 dimension vectors are distinct. The seven further indecomposables
with central dimension zero consist of the simple at $J$ and the three
interval representations on each two-vertex arm. Together these give
36 distinct indecomposables over an algebraic closure. Gabriel's theorem
and the 36 positive roots of $E_6$ establish completeness. The fixed root
reflection check in production is an additional consistency check; no
input-dependent root or representation classification remains to be solved.
The cited representation-theory text states Gabriel's theorem at
Theorem 6.5.2 and gives the $E_6$ positive-root count immediately before it
\cite{etingof}. The multiple-flag classification provides a separate classical
antecedent for the $(2,3,3)$ case \cite{mwz}; the classification itself is not
new research of this continuation.

The finite-field descent is valid. Write $K=\F_q$, and let $V$ be an
arbitrary representation over $K$. Over its algebraic closure, $V$ is
isomorphic to a sum $M$ of the displayed prime-field representatives. Only
finitely many algebraic coefficients occur in this isomorphism and its
inverse, so the isomorphism is defined over some finite extension $E/K$,
of degree $r$. Restricting scalars gives
\[
 V^{\oplus r}\simeq M^{\oplus r}
 \quad\text{as representations over }K.
\]
Krull--Schmidt over $K$ equates the multiplicities of every indecomposable
on both sides, hence gives $V\simeq M$. The needed arbitrary-field and
descent statements are precisely the material in Problems 3.8.3--3.8.4
of \cite{etingof}. There are therefore no extra twisted types over
$\F_{2^a}$ and no hidden enumeration of field elements to descend the
classification. Injectivity of each arrow is inherited by a direct summand,
so the seven zero-central types cannot occur in the subspace problem.

\section{Phase orientation and the local signed polynomial}
The inherited phase acquisition deserves particular attention because a
transpose error would survive many symmetric examples. In the recovered
proof, if the multiplier is $M$ and the character is
$(-1)^{\operatorname{Tr}\operatorname{tr}(BM)}$, the geometric variable is
$A=M^{\mathsf T}$ and the character matrix is $Q=B^{\mathsf T}$.
The target contribution is $CA$. Left translates of the character have
annihilator $MB=0$, which becomes $QA=0$. Their image carrier is therefore
\[
 F=C\ker Q,
\]
with precisely this orientation. The production inherited
\path{acquire_matrix_phase_carriers} uses actual left products in the
multiplier algebra and the acquired binary character. It does not assume
that $Q$ is symmetric or replace its kernel by that of its transpose.

For a candidate $J\le E_1+E_2$, set
$\ell_i=\dim(J\cap E_i)$ and $\eta_i=\dim(J\cap F_i)$.
For one factor, abbreviate $s=s_i$, $g=\rank C$, $t=\rank Q$,
$f=\dim C\ker Q$. The columns of an allowed matrix $A$ lie in
\[
 L_J=C^{-1}(J\cap E),\qquad \dim L_J=s-g+\ell.
\]
The map $C|_{\ker Q}$ has kernel dimension $s-t-f$. Consequently
\[
 \dim(L_J\cap\ker Q)=s-t-f+\eta,
 \qquad
 \rank(Q|_{L_J})=t-g+\ell+f-\eta.
\]
This proves the claimed restricted phase rank and matches the production
assignment to \path{tau}. No equality of $E_1$ and $E_2$ is used.

Choose an $r$-column basis for $L_J$, with $r=s-g+\ell$. In these
coordinates $A$ is an arbitrary $r$-by-$s$ matrix, and the restricted
character has rank $\tau$. Summing a character over matrices with image
inside a fixed $v$-space yields $q^{sv}$ precisely when that space lies
in the character kernel, and zero otherwise. Summing this statement over
all $v$-spaces gives
\[
 \gb{r-\tau}{v}q^{sv}
   =\sum_{u=0}^{v}P_u\gb{r-u}{v-u}.
\]
The triangular system has diagonal one. Thus the reused local routine
computes exactly the signed rank polynomial, and multiplication of the
two local polynomials is justified by independence of the two direct
algebra factors. Zero target carriers, zero phase, full phase rank,
$r=0$, and a zero ambient join are covered by these formulas.

\section{Acquisition of actual geometry, not a dimension surrogate}
The native route first acquires the nucleus, the affine Hom slice and its
trace-annihilator, the multiplier algebra, its two central factors, their
actual common scalar action, regular component dimensions and zero common
kernel. The one-copy orbit test is performed on actual vectors obtained
from $Dy$ and nucleus operators. The phase carrier is constructed from
the actual translated-character annihilator. These operations create
$F_i\subseteq E_i\subseteq S^*$ as right $S$-submodules.

Only after that module origin is established is a binary dimension divided
by $ap$. Morita equivalence sends these submodules to ordinary finite-field
subspaces, with natural dimension equal to that quotient. Their lattice
has field order $q$, not $2^{ap}$. This distinction is correctly preserved
by the code. The lower-level histogram helper consumes certified module
data. Its simple divisibility checks are not a standalone recognition
theorem for arbitrary binary subspaces; this is not a gap in the native
wrapper because the wrapper manufactures the modules.

For the two flags $0\subseteq F_i\subseteq E_i\subseteq R$, their
relative position is fully described by the nine mixed differences
\[
 b_{uv}=\dim(V_u\cap W_v)-\dim(V_{u-1}\cap W_v)
       -\dim(V_u\cap W_{v-1})+\dim(V_{u-1}\cap W_{v-1}).
\]
Successive complements in the intersections give a simultaneous basis
with these multiplicities. This proves nonnegativity and completeness of
the invariant for a pair of flags. Production computes the intersections
from actual binary row spaces, checks nesting and integral dimensions,
and uses their natural dimensions. Since $R=E_1+E_2$, the outside-both
cell $b_{33}$ vanishes. It is not inferred from the individual ranks of
the four subspaces alone.

\section{Automorphism weights and the exact census}
The implementation of \path{hom_dimension} solves the commuting-square
equations $f_vA=Bf_u$ for every arrow, with separate variables for each
vertex map. Its row and column indexing matches these equations. Since
all fixed inclusion matrices have entries in $\F_2$, the binary rank
of each coefficient matrix is unchanged by extension to $\F_{2^a}$.
The resulting fixed $29$-by-$29$ Hom table is therefore valid for all
admitted fields.

For a multiplicity vector $m$, let $M=\bigoplus_\alpha T_\alpha^{m_\alpha}$.
Its endomorphism dimension is
$\sum_{\alpha,\beta}m_\alpha m_\beta h_{\alpha\beta}$.
Because distinct listed types are nonisomorphic and each has endomorphism
ring exactly $\F_q$, the semisimple quotient of $\End(M)$ is
$\prod_\alpha M_{m_\alpha}(\F_q)$. The radical has the dimension
used in production. Lifting units through this nilpotent radical gives
\[
 |\Aut(M)|=
 q^{\sum m_\alpha m_\beta h_{\alpha\beta}-\sum m_\alpha^2}
 \prod_\alpha |\GL_{m_\alpha}(q)|.
\]
This is an ordinary automorphism-group calculation, not a supplied
Hall polynomial.

Let $G$ be the automorphism group of the actual pair of base flags.
The same group is the automorphism group of their representation with
$J=0$, so its order is obtained from the thin-only multiplicities.
Every allowed full multiplicity vector gives one isomorphism class of
the complete configuration. Since the base flags are fixed, this is
one $G$-orbit on $J$. Conversely, two $J$'s in the same class are
related by an ambient automorphism preserving both flags; all other
vertex maps are determined by the ambient map. Hence the stabilizer
is exactly $\Aut(M)$, and
\[
 \#\{J\text{ in that orbit}\}=|G|/|\Aut(M)|.
\]
There is no missing basis-choice factor and no extra quotient by
$\Aut(J)$. The injective inclusions already force the induced map on $J$.
Production checks exact division and sums this positive weight at the
five additive intersection statistics. Its Gaussian totals for each
$\dim J$ are valuable consistency checks, although those totals alone
would not prove the finer census.

The recursive enumeration is also faithful: all eleven non-thin
multiplicities are chosen, leaving a nonnegative grid which is split
between the two thin types in each cell. Every resulting multiplicity
vector is emitted exactly once. All entries of the mutable profile are
assigned before each emitted leaf; no stale choice survives from a
previous recursion branch.

\section{Join inversion and the original-input answer}
The census yields moments over all containing $J$ of dimension $j$.
A fixed actual image of dimension $\rho$ occurs in exactly
$\gb{k-\rho}{j-\rho}$ such spaces. Thus the triangular inversion
\[
 H_j=M_j-\sum_{\rho<j}\gb{k-\rho}{j-\rho}H_\rho
\]
is valid coefficientwise even when the coefficients are signed. The
production loop uses this precise scale. The exponent of its local
rank polynomial is $h_i\rank A_i$, giving the natural rank of the
regular dual operator; its second statistic is the natural rank of the
actual joined target image. These are the two ranks required by the
retained signed universal-product fiber.

The native call uses the acquired affine Hom offset, target constants,
and output-dependency relations. The corrected affine-coordinate
reconstruction from the retained frontend is therefore the relevant
one; the obsolete polar-coordinate offset is not used. No completed Hom
coordinate and no gate wire becomes an extra independent input.
The production result divides by $2^c$, checks integrality and the
original-input range $[0,2^n]$, and returns the complement using the
same original $n$.

If the target violates output dependencies, the early zero return is
sound even before the exactly-two-factor consumer is entered. It should
be interpreted as an independent zero certificate, not as evidence that
an arbitrary unsupported factor configuration acquired an $E_6$ model.

\section{Uniform cost and exclusions}
There are only eleven non-thin types, and every multiplicity is bounded
by $k$. The remaining thin splits use at most eight occupied cells because
$b_{33}=0$. Thus the claimed $(k+1)^{19}$ profile bound is a genuine fixed
polynomial. The census uses at most $(k+1)^5$ statistic positions, and
the later polynomial degrees and moment indices are bounded by $d$ and
$k$. None of these bounds contains $q$ as an enumeration range.

Automorphism integers have $O(ak^2)$ bits. All signed matrix sums are
bounded in absolute value by the number of matrices contributing to
them, and the containing-space factors add only $O(ak^2)$ bits.
The stated $O(a(k^2+s_1^2+s_2^2)+n)$ intermediate bound, and the looser
$O(N^2)$ bound used in the theorem, are adequate. Gaussian coefficients,
group orders, exact quotients and signed cancellation all have polynomial
bit cost. Combined with the fixed profile exponent and retained algebra
acquisition, the conservative $O(N^{30})$ envelope is sufficient; it is
not an optimized or empirical runtime estimate.

The matrix-algebra certification here uses actual multiplication,
Frobenius and sandwich-span linear algebra. It does not manufacture
matrix units or search for an invertible dual generator. Consequently
the random primitive-idempotent manufacture used by the separate
constructive section is not silently imported into this deterministic
counting theorem. The proof's existential $D_0$ is sufficient for its
rank argument because every production quantity is acquired without it.

The exclusions remain substantive:
\begin{itemize}
\item Exactly two regular matrix-dual factors with the same actual scalar
field as the nucleus are required. The older varied-field commutative
semisimple branch cannot be substituted merely by matching dimensions.
\item Each target component is zero or occupies one joint simple copy.
Targets spread across several copies are not covered by this theorem.
\item Positive natural multiplicities $h_i$ suffice here; the condition
$h_i\ge2$ belongs to the separate affine-section construction and should
not be added to this count theorem.
\item Ordinary-circuit syntax and the acquired affine Hom conditions are
retained admission conditions. Two global split orientations are tried;
a legitimate decline is possible outside them.
\item This construction does not give a uniform sampler or arbitrary
original-input conditional counts. Signed rank coefficients are not
sampling weights, and arbitrary conditioning need not preserve the dual.
\item Three general factors and unrestricted Boolean circuits remain
outside the proved polynomial scope. The finite $E_6$ catalogue cannot
be reused for a different quiver without a governing argument.
\end{itemize}

\section{Evidence read and its limits}
The saved audit reports all assertions passed. Its geometry oracle
enumerates physical finite-field subspaces independently of the quiver
census, covering 825 canonical pair-flag types over $\F_2,\F_4,\F_8$,
8,015 census coefficients and 47,949 actual subspace instances.
The separate matrix oracle covers 36 cases and 104,832 local matrices,
joining actual column spaces and evaluating phases. Those cases include
the ambient-three types and nonsymmetric phase matrices, making them
relevant to classification and orientation rather than only final mass.

The native fixture has 72 original inputs and 24,635 gates. Its target
carriers are distinct lines; the old diamond consumer declines with
``nonzero target carriers do not coincide''. The new count is
\[
 71\,228\,562\,270\,584\,832.
\]
The separate original-left-matrix/right-linear-system oracle agrees.
The saved audit also checks all 256 affine phase words at this target,
all 256 pairs of binary $2$-by-$2$ target matrices with independently
varied phases, 128 random original-gate evaluations and 32 constructed
positive original-gate evaluations. The audit harness's frontier
permutation/constant correction is retained in
\path{new_work/phase_constructive/audit_failures.md}. The positive
checks matter because a low-density predicate can make exclusively
random tests weak.

The saved formula-growth case uses $q=2^{17}$, two size-six factors,
codimension 1,224, and an 8,568-bit answer. It is explicitly marked
formula-only: the family parameter 11,016 original bits is not an
executed raw-DAG acquisition. No independent broad audit rerun was
performed as part of this review. These are correctness records, not
evidence of practical superiority: in particular the native fixture's
codimension-eight character fallback is already feasible. The barrier
removed by the theorem is the target-carrier equality restriction at
growing structural parameters, with recognition and bit cost retained.

\section*{Retained files inspected}
In addition to the four hashed files above, this review inspected the
recovered production files
\path{work_matrix_root/matrix_affine_flag_counter.py},
\path{work_matrix_root/matrix_target_geometry.py}, and
\path{work_matrix_root/finite_simple_binary_algebra.py}, and their
governing affine-flag and target-geometry arguments. The recovered
regular-module, affine-input and universal-product results remain the
explicit dependencies of the new theorem. Nothing in this disposition
converts earlier progress-record reports lacking governing arguments
into established results.

\begin{thebibliography}{2}
\bibitem{etingof} P. Etingof et al., \emph{Introduction to Representation
Theory}, Theorem 6.5.2, Example 6.4.9, Problems 3.8.3--3.8.4.
\url{https://math.mit.edu/~etingof/representationtheorybook.pdf}.
Primary text checked on 28 October 2026.
\bibitem{mwz} P. Magyar, J. Weyman, A. Zelevinsky,
\emph{Multiple Flag Varieties of Finite Type}, Adv. Math. 141 (1999),
97--118, Theorems 2.4, 2.6 and 2.9, especially the $(2,3,3)$ case.
\url{https://arxiv.org/abs/math/9805067}.
Primary preprint checked on 28 October 2026.
\end{thebibliography}
\end{document}
